% ============================================================
%  IAM — A Note on Entropic Gravity and the
%         Thermodynamic-Gravity Conjecture
%  Heath W. Mahaffey, Independent Researcher, Entiat WA
%  Prepared for Prof. E. N. Saridakis — March 2026
% ============================================================

\documentclass[12pt]{article}

% ── Packages ────────────────────────────────────────────────
\usepackage[margin=1in]{geometry}
\usepackage{amsmath, amssymb, amsthm}
\usepackage{graphicx}
\usepackage{booktabs}
\usepackage{array}
\usepackage{xcolor}
\usepackage{hyperref}
\usepackage{microtype}
\usepackage{parskip}
\usepackage{setspace}
\usepackage{titlesec}
\usepackage{fancyhdr}
\usepackage{natbib}

% ── Colors ──────────────────────────────────────────────────
\definecolor{iamblack}{RGB}{46,74,122}
\definecolor{lightgray}{RGB}{245,247,250}
\definecolor{midgray}{RGB}{120,120,120}

% ── Hyperref setup ──────────────────────────────────────────
\hypersetup{
  colorlinks = true,
  linkcolor  = iamblue,
  citecolor  = iamblue,
  urlcolor   = iamblue,
  pdftitle   = {IAM: A Note on Entropic Gravity and the Thermodynamic-Gravity Conjecture},
  pdfauthor  = {Heath W. Mahaffey}
}

% ── Section formatting ──────────────────────────────────────
\titleformat{\section}
  {\large\bfseries\color{iamblue}}
  {\thesection.}{0.6em}{}[\vspace{-4pt}\rule{\linewidth}{0.4pt}]

\titleformat{\subsection}
  {\normalsize\bfseries\color{iamblue}}
  {\thesubsection.}{0.5em}{}

% ── Header / footer ─────────────────────────────────────────
\pagestyle{fancy}
\fancyhf{}
\fancyhead[L]{\small\color{midgray}Mahaffey — IAM: Entropic Gravity Note}
\fancyhead[R]{\small\color{midgray}March 2026}
\fancyfoot[C]{\small\thepage}
\renewcommand{\headrulewidth}{0.4pt}
\renewcommand{\headrule}{\hbox to\headwidth{\color{iamblue}\leaders\hrule height \headrulewidth\hfill}}

% ── Spacing ─────────────────────────────────────────────────
\setlength{\parskip}{6pt}
\setlength{\parindent}{0pt}
\onehalfspacing

% ── Math shortcuts ──────────────────────────────────────────
\newcommand{\Oms}{\Omega_m}
\newcommand{\bm}{\beta_m}
\newcommand{\Ea}{E(a)}
\newcommand{\Sbh}{S_\mathrm{BH}}
\newcommand{\Sinfo}{S_\mathrm{info}}
\newcommand{\Stot}{S_\mathrm{total}}
\newcommand{\THawk}{T_H}
\newcommand{\AHawk}{A_H}
\newcommand{\LCDM}{$\Lambda$CDM}
\newcommand{\musig}{$\mu$--$\Sigma$}

% ── Theorem-like environments ───────────────────────────────
\newtheorem{result}{Result}

% ============================================================
\begin{document}

% ── Title block ─────────────────────────────────────────────
\begin{center}
  {\LARGE\bfseries\color{iamblue}
    Informational Actualization Model}\\[6pt]
  {\large\itshape
    A Note on Entropic Gravity and the\\
    Thermodynamic-Gravity Conjecture}\\[10pt]
  {\normalsize
    Heath W.\ Mahaffey\\
    Independent Researcher, Entiat, WA\\
    \href{mailto:hmahaffeyges@gmail.com}{\texttt{hmahaffeyges@gmail.com}}}\\[4pt]
  {\small\color{midgray}
    Zenodo: \href{https://doi.org/10.5281/zenodo.18702042}{\texttt{10.5281/zenodo.18702042}}
    \quad|\quad
    GitHub: \href{https://github.com/hmahaffeyges/IAM-Validation}{\texttt{hmahaffeyges/IAM-Validation}}
    \quad|\quad
    March 2026}
\end{center}

\vspace{4pt}
{\color{iamblue}\rule{\linewidth}{1pt}}
\vspace{4pt}

% ── Abstract ────────────────────────────────────────────────
\begin{center}
  \small
    \textbf{Abstract}\\
\end{center}
\begin{center}
  \begin{minipage}{0.80\linewidth}
    \small
    \textbf\\
    The Informational Actualization Model (IAM) is a cosmological extension of the
    Jacobson--Cai-Kim thermodynamic gravity framework.  Rather than modifying the
    entropy functional $\Sbh$ on the cosmic horizon --- the approach taken in Barrow
    and Tsallis holographic cosmologies --- IAM identifies a physically motivated
    additional entropy source, $\Sinfo$, arising from gravitational decoherence of
    collapsing matter.  This entropy enters the Clausius relation and propagates to the
    perturbation equations as a single Hubble friction term with no new free parameters:
    the coupling constant $\bm = \Oms/2$ is derived from the virial theorem, and the
    activation function $\Ea = \exp(1-1/a)$ from the decoherence integral over cosmic
    history.  This note situates IAM within the entropy-gravity framework familiar from
    holographic dark energy research and identifies the key structural difference between
    entropy-\emph{modification} approaches and entropy-\emph{source} approaches.
  \end{minipage}
\end{center}

\vspace{4pt}
{\color{iamblue}\rule{\linewidth}{0.6pt}}
\vspace{8pt}

% ============================================================
\section{The Shared Starting Point}

IAM begins from the same foundation as holographic dark energy models in the
Barrow--Tsallis tradition: the \emph{thermodynamic-gravity conjecture}.  In this
framework the standard Friedmann equations governing cosmic expansion are not
fundamental --- they are thermodynamic equations of state derivable from the first
law of thermodynamics applied to the cosmic horizon.

The derivation proceeds through the Cai--Kim (2005) formalism.  The apparent horizon
of a Friedmann--Robertson--Walker universe has area $\AHawk = 4\pi/H^2$,
Gibbons--Hawking temperature $\THawk = H/2\pi$, and Bekenstein--Hawking entropy
$\Sbh = \AHawk/4$.  Applying the first law
\begin{equation}
  \mathrm{d}E = \THawk\,\mathrm{d}\Sbh - W\,\mathrm{d}V
  \label{eq:first_law}
\end{equation}
recovers the standard Friedmann equations exactly.  Gravitational field equations
are thermodynamic identities \citep{Jacobson1995}.

This starting point is shared by IAM, Barrow holographic dark energy, Tsallis
holographic dark energy, and Kaniadakis cosmology.  All of these frameworks accept
the thermodynamic-gravity conjecture.  They diverge in what they do next.

% ============================================================
\section{Two Roads From the Conjecture}

Once the thermodynamic-gravity conjecture is accepted there are two natural ways to
extend it beyond \LCDM:

\subsection{Road 1 --- Modify the entropy functional}

Barrow entropy replaces $\Sbh$ with a fractal-deformed expression
$S_\mathrm{Barrow} = (A/A_0)^{1+\Delta/2}$, introducing a deformation parameter
$\Delta$ that quantifies quantum-gravitational fractal structure on the horizon
surface \citep{Saridakis2020barrow}.  Tsallis entropy replaces $\Sbh$ with a
non-extensive expression $S_\mathrm{Tsallis} = \gamma A^\delta$, where $\delta$
captures long-range correlations \citep{Saridakis2018tsallis}.  In both cases the
modification is to the geometric entropy term itself, and the new parameter
($\Delta$ or $\delta$) must be constrained by data.

These are well-motivated approaches.  They explore what happens when the
entropy-area law receives corrections --- a natural question if the horizon geometry
is not smooth at the Planck scale.

\subsection{Road 2 --- Identify an additional entropy source}

IAM takes a different path.  The geometric entropy $\Sbh$ remains standard
Bekenstein--Hawking.  Instead, IAM identifies a physically motivated additional
contribution to the total horizon entropy budget:
\begin{equation}
  \Stot = \Sbh + \Sinfo,
  \label{eq:stotal}
\end{equation}
where $\Sinfo$ arises from \emph{gravitational decoherence} of collapsing matter.

The reasoning is as follows.  Collapsing matter continuously generates classical
information: quantum states decohere into classical configurations as gravitational
gradients cause wave-function branching.  Each decoherence event is irreversible
and carries an entropy cost of $k_B \ln 2$ per bit (Landauer's principle,
experimentally verified by \citealt{Berut2012} and \citealt{Jun2014}).  This
information is holographically encoded on the cosmic horizon at the
Gibbons--Hawking temperature $\THawk = H/2\pi$.

$\Sinfo$ is not a modification to the entropy-area law.  It is a \emph{source
term} --- entropy entering the horizon accounting from a physical process that
standard thermodynamic gravity had not previously included.

% ============================================================
\section{Where the Two Roads Lead}

The distinction between Road~1 and Road~2 is not merely taxonomic --- it has
direct observational consequences.

In Barrow and Tsallis cosmologies the entropy modification alters the Friedmann
equations --- the background expansion history is modified.  The new parameter
($\Delta$ or $\delta$) affects both the expansion rate $H(z)$ and the growth of
structure simultaneously, and must be fitted to data.  Recent constraints from
DESI DR2 BAO data confirm that both Barrow and Tsallis cosmologies are compatible
with observations but find that neither resolves the $H_0$ tension, and that both
are disfavoured relative to \LCDM\ by information criteria due to the additional
free parameter \citep{Luciano2025}.

In IAM, $\Sinfo$ is sourced by matter decoherence and therefore enters the
perturbation equations only --- specifically as an additional Hubble friction term
in the matter growth equation:
\begin{equation}
  \delta'' + \bigl[2 + \bm\,\Ea\bigr]\,H\,\delta' - \tfrac{3}{2}\,\Oms H^2\,\delta = 0,
  \label{eq:growth}
\end{equation}
where primes denote derivatives with respect to cosmic time.  The background
expansion remains standard \LCDM.  The photon sector is unaffected.

This is the origin of the dual-sector \musig\ prediction.  Massive particles
travel on timelike paths; proper time elapses; gravitational decoherence acts
continuously.  Photons travel on null geodesics; no proper time elapses; they are
never subject to the decoherence process.  The prediction $\Sigma = 1$ (no
lensing-potential modification) follows from this geometric fact, not from a
parameter choice.  The prediction $\mu < 1$ (suppressed matter growth) follows
from the friction term in Eq.~\eqref{eq:growth}.

% ============================================================
\section{The Coupling Constant \texorpdfstring{$\beta_m = \Omega_m/2$}{beta\_m}}

The coefficient $\bm$ in Eq.~\eqref{eq:growth} is not a free parameter.  It is
derived from the virial theorem.

For any self-gravitating system bound by a $1/r$ potential:
\begin{equation}
  2\langle T\rangle + \langle V\rangle = 0
  \implies
  \langle T\rangle = \tfrac{1}{2}|\langle V\rangle|.
  \label{eq:virial}
\end{equation}
Exactly half the gravitational energy budget resides in kinetic degrees of freedom.
In IAM the kinetic channel is the decoherence-producing channel --- irreversible
random motions of particles within collapsing structures generate the classical
information requiring holographic encoding.  The informational energy density
contribution is therefore $\Oms/2$, giving $\bm = 0.1575$.

Six independent N-body simulation groups have measured the mass-weighted virial
ratio $2K/|U|$ directly \citep{Bryan1998, Bett2007, Neto2007, Ludlow2010,
Power2012, Klypin2016}, with combined result $\langle\eta_\mathrm{virial}\rangle =
0.815\pm0.025$, giving $\bm = 0.159\pm0.010$.  This agrees with the
virial-theorem-derived $\Oms/2 = 0.1575$ to $1.1\%$.  The same $1/2$ partition
holds from hydrogen atoms through stellar structure through galaxy clusters ---
37~orders of magnitude in physical scale.

The Planck~2018 MCMC posterior independently recovers $\bm = \Oms/2$ to within
$0.2\sigma$ without this value being fitted.  That independent confirmation from
the most constraining cosmological dataset is the single most compelling
quantitative result in the framework.

% ============================================================
\section{The Activation Function \texorpdfstring{$E(a)$}{E(a)}}

$\Ea$ is derived by integrating the decoherence rate over cosmic history, weighted
by the Gibbons--Hawking encoding temperature $\THawk = H/2\pi$ and horizon area
$\AHawk = 4\pi/H^2$.  The result is:
\begin{equation}
  E(a) = \exp\!\left(1 - \frac{1}{a}\right).
  \label{eq:Ea}
\end{equation}

Two boundary conditions fix $\Ea$ completely without free parameters:
\begin{itemize}
  \item $E(1) = 1$: the present-day informational contribution is normalized to
        the matter energy density scale.
  \item $E(a) \to e$ as $a\to\infty$: the self-limiting nature of the feedback
        (expansion dilutes matter, slows decoherence) causes cumulative entropy
        to asymptote to a finite ceiling.
\end{itemize}

At $a=1$ today, $E(1)/e = 1/e \approx 36.8\%$, placing the current epoch at the
\emph{inflection point} of the activation function where the rate of informational
entropy production is at its peak.  Six published structure-formation analyses find
an effective growth exponent $n_\mathrm{eff} = 3.22\pm0.44$, consistent with the
range $3.0$--$3.5$ required by the decoherence integral.

% ============================================================
\section{Comparison With Barrow/Tsallis Holographic Cosmologies}

Table~\ref{tab:compare} summarizes the structural comparison.  Both families of
models emerge from the same thermodynamic-gravity conjecture.  The key difference
is where new physics enters: in the entropy functional (Barrow/Tsallis) versus in
the entropy source term (IAM).

\begin{table}[h]
  \centering
  \caption{Structural comparison of entropy-based cosmological frameworks.}
  \label{tab:compare}
  \renewcommand{\arraystretch}{1.35}
  \begin{tabular}{>{\raggedright}p{3.2cm}
                  >{\raggedright}p{4.8cm}
                  >{\raggedright\arraybackslash}p{5.0cm}}
    \toprule
    \textbf{Feature}
      & \textbf{Barrow / Tsallis HDE}
      & \textbf{IAM} \\
    \midrule
    Foundation
      & Thermodynamic-gravity conjecture (Jacobson; Cai--Kim)
      & Same foundation \\
    Modification
      & Entropy functional: $\Sbh \to S_\mathrm{Barrow}$ or $S_\mathrm{Tsallis}$
      & Entropy source: $\Stot = \Sbh + \Sinfo$ \\
    Free parameters
      & $\Delta$ (Barrow) or $\delta$ (Tsallis) --- fitted to data
      & None --- $\bm = \Oms/2$ from virial theorem \\
    Background modified?
      & Yes --- Friedmann equations altered
      & No --- standard \LCDM\ background \\
    Perturbations modified?
      & Yes (through modified Friedmann equations)
      & Yes --- direct friction term only \\
    \musig\ signature
      & Generally $\mu\ne1$, $\Sigma\ne1$
      & $\mu < 1$,\ $\Sigma = 1$ exactly \\
    $H_0$ tension
      & Not resolved \citep{Luciano2025}
      & Predicted: $72.5$~km/s/Mpc (output) \\
    AIC vs.\ \LCDM
      & Disfavoured (extra parameter)
      & Indifferent (no extra parameter) \\
    \bottomrule
  \end{tabular}
\end{table}

The AIC comparison deserves particular attention.  Because Barrow and Tsallis
cosmologies introduce one additional free parameter, information criteria penalise
them relative to \LCDM\ even when the fit improves.  IAM introduces no additional
free parameter --- $\bm$ is fully determined by $\Oms$, which is already a
standard \LCDM\ parameter.  IAM therefore carries no AIC penalty.

% ============================================================
\section{A Note on the Entropy Exponent}

There is a natural question to ask of the Barrow/Tsallis framework: why does the
entropy exponent take the value it does?  The deformation parameters $\Delta$ and
$\delta$ are physically motivated --- fractal horizon structure, non-extensive
statistics --- but their values are not derived from first principles.  They are
measured from data.

IAM suggests a perspective on an adjacent question.  If $\Sinfo$ from decoherence
is the correct physical source of the cosmological entropy correction, then the
apparent need for a non-standard entropy functional in other models may reflect an
incomplete accounting of entropy \emph{sources} rather than a modification to the
entropy-area law itself.  Standard Bekenstein--Hawking entropy may be exact; the
horizon entropy budget simply includes more terms than previously recognised.

Whether Barrow or Tsallis entropy captures genuine Planck-scale physics is an
open question IAM does not resolve.  The narrower point is that the same late-time
cosmological phenomenology that motivates entropy-modification approaches can be
obtained from an entropy-source approach with zero free parameters and a derivable
coupling constant.

% ============================================================
\section{MCMC Validation and Key Predictions}

\subsection{Validation}

IAM has been validated against the Planck~2018 CMB likelihood through 17 converged
MCMC chains (12 Level~1 via MGCAMB \musig\ parametrisation; 3 Level~2 via direct
Fortran modification of CAMB's perturbation equations; 2 exploratory background
runs confirming the perturbation-level mechanism).  All chains satisfy
$R-1 < 0.01$ Gelman--Rubin convergence.

\begin{table}[h]
  \centering
  \caption{Key IAM predictions and current observational status.}
  \label{tab:numbers}
  \renewcommand{\arraystretch}{1.35}
  \begin{tabular}{lcll}
    \toprule
    \textbf{Quantity} & \textbf{IAM Value} & \textbf{Current Data} & \textbf{Source} \\
    \midrule
    $\mu_0$          & $-0.136$ (derived)     & $0.05\pm0.22$            & DESI 2025 \\
    $\Sigma_0$ (dev.)& $0$ (null worldlines)  & $0.008\pm0.045$          & ACT+WMAP+SDSS \\
    $\sigma_8$       & $0.800$ (derived)      & $0.802\pm0.020$          & St\"{o}lzner et al.\ 2025 \\
    $H_0^\mathrm{matter}$ & $72.5$~km/s/Mpc (derived) & $73.04\pm1.04$    & Riess et al.\ 2022 \\
    $\Delta\chi^2$   & $+0.54$                & confirmed                & 17 Planck chains \\
    $\bm$            & $\Oms/2 = 0.1575$      & $0.159\pm0.010$          & 6 N-body studies \\
    \bottomrule
  \end{tabular}
\end{table}

\subsection{Falsifiable predictions}

IAM makes three near-term quantitative predictions:

\begin{enumerate}
  \item \textbf{Euclid \musig\ measurement.}  The IAM prediction
        $\mu_0 = -0.136$, $\Sigma_0 = 0$ is directly testable by Euclid at
        projected $\sigma(\mu_0)\sim0.04$ sensitivity.  The combination
        $\mu < 1$ with $\Sigma = 1$ is a unique signature in the \musig\ plane:
        $f(R)$ gravity gives $\mu > 1$, $\Sigma > 1$; DGP braneworld gives
        $\mu > 1$; general Horndeski theories predict $\Sigma\ne1$.

  \item \textbf{DESI DR2 growth rate trajectory.}  The IAM $f\sigma_8(z)$
        trajectory, suppressed by the $\bm\Ea$ friction term, is consistent
        with DESI DR2 growth rate data without requiring $w < -1$.  A direct
        confrontation of the DESI DR2 $f\sigma_8$ measurements against the
        IAM predicted trajectory tests the interpretation that the apparent
        phantom crossing at $z\sim0.5$ is a growth-geometry tension, not a
        dark energy signature.

  \item \textbf{Galaxy cluster lensing-to-dynamics ratio.}  The dual-sector
        structure predicts $M_\mathrm{lens}/M_\mathrm{dyn} = 1/\mu_\mathrm{eff}(z) > 1$
        in a redshift-dependent pattern that approaches unity at high $z$ as
        $\mu(a)\to1$.  This is qualitatively distinct from a constant
        hydrostatic bias and testable with current eROSITA, Planck~SZ, and
        weak lensing datasets in combination.
\end{enumerate}

If Euclid measures $\mu$ consistent with zero at $\sigma(\mu_0)\sim0.04$
precision, IAM is ruled out.  The author is not aware of any currently published modified gravity model
that predicts $\mu < 1$ with $\Sigma = 1$ from first principles.

% ============================================================
\section{Summary}

The shared foundation is the Jacobson--Cai-Kim result: gravitational field equations are
thermodynamic equations of state derivable from the first law applied to the cosmic
horizon.

IAM's specific contribution is to identify gravitational decoherence of collapsing
matter as a physical \emph{source} of additional horizon entropy.  This source
term, derived via Landauer's principle and the virial theorem, produces a
matter-sector Hubble friction term with no free parameters.  The coupling
$\bm = \Oms/2$ is recovered independently by six N-body simulation groups and by
the Planck~2018 MCMC posterior to $0.2\sigma$.

The framework is consistent with the full Planck~2018 likelihood
($\Delta\chi^2 = +0.54$), predicts $\sigma_8 = 0.800$ in agreement with current
weak lensing measurements, and places the $H_0$ tension endpoints within $1\sigma$
of both Planck and SH0ES with no parameter adjustment.  It carries no AIC penalty
because no new free parameter is introduced.

Euclid's \musig\ measurement at projected $\sigma(\mu_0)\sim0.04$ sensitivity
provides a direct test of the unique IAM prediction $\mu_0 = -0.136$,
$\Sigma_0 = 0$.  The full derivation chain, MCMC chains, modified CAMB Fortran
source, and analysis scripts are publicly available at the Zenodo DOI and GitHub
repository listed on the title page.

% ============================================================
\section*{Acknowledgements}

The author thanks Prof.\ Emmanuel N.\ Saridakis, whose development of Barrow
holographic dark energy and the broader entropy-based cosmological programme
demonstrated that the thermodynamic-gravity conjecture is a productive and
observationally viable foundation for extensions beyond \LCDM.  The clarity
and rigour of that body of work made the structural question addressed here ---
whether new cosmological physics enters through the entropy functional or through
an unaccounted entropy source --- tractable in a way it would not otherwise have been.

% ============================================================
\bibliographystyle{plainnat}
\begin{thebibliography}{99}

\bibitem[Bérut et al.(2012)]{Berut2012}
B\'{e}rut, A., et al., 2012,
\textit{Experimental verification of Landauer's principle linking information and thermodynamics},
Nature, 483, 187.

\bibitem[Bett et al.(2007)]{Bett2007}
Bett, P., et al., 2007,
\textit{The spin and shape of dark matter haloes in the Millennium simulation},
MNRAS, 376, 215.

\bibitem[Bryan \& Norman(1998)]{Bryan1998}
Bryan, G.~L.\ \& Norman, M.~L., 1998,
\textit{Statistical properties of X-ray clusters: analytic and numerical comparisons},
ApJ, 495, 80.

\bibitem[Cai \& Kim(2005)]{CaiKim2005}
Cai, R.-G.\ \& Kim, S.~P., 2005,
\textit{First law of thermodynamics and Friedmann equations of FRW universe},
JHEP, 02, 050.

\bibitem[Jacobson(1995)]{Jacobson1995}
Jacobson, T., 1995,
\textit{Thermodynamics of spacetime: the Einstein equation of state},
Phys.\ Rev.\ Lett., 75, 1260.

\bibitem[Jun et al.(2014)]{Jun2014}
Jun, Y., Gavrilov, M., \& Bechhoefer, J., 2014,
\textit{High-precision test of Landauer's principle in a feedback trap},
Phys.\ Rev.\ Lett., 113, 190601.

\bibitem[Klypin et al.(2016)]{Klypin2016}
Klypin, A., et al., 2016,
\textit{MultiDark simulations: the story of dark matter halo concentrations and density profiles},
MNRAS, 457, 4340.

\bibitem[Luciano et al.(2025)]{Luciano2025}
Luciano, G.~G., Paliathanasis, A., \& Saridakis, E.~N., 2025,
\textit{Barrow and Tsallis entropies after the DESI DR2 BAO data},
JCAP, 09, 013.

\bibitem[Ludlow et al.(2010)]{Ludlow2010}
Ludlow, A.~D., et al., 2010,
\textit{The formation of \LCDM\ haloes from the inside out},
MNRAS, 406, 137.

\bibitem[Neto et al.(2007)]{Neto2007}
Neto, A.~F., et al., 2007,
\textit{The statistics of \LCDM\ halo concentrations},
MNRAS, 381, 1450.

\bibitem[Planck Collaboration(2020)]{Planck2020}
Planck Collaboration, 2020,
\textit{Planck 2018 results VI: cosmological parameters},
A\&A, 641, A6.

\bibitem[Power et al.(2012)]{Power2012}
Power, C., et al., 2012,
\textit{Spurious evolution in cosmological N-body simulations},
MNRAS, 419, 1576.

\bibitem[Riess et al.(2022)]{Riess2022}
Riess, A.~G., et al., 2022,
\textit{A comprehensive measurement of the local value of the Hubble constant},
ApJL, 934, L7.

\bibitem[Saridakis(2020a)]{Saridakis2020barrow}
Saridakis, E.~N., 2020,
\textit{Barrow holographic dark energy},
Phys.\ Rev.\ D, 102, 123525.

\bibitem[Saridakis et al.(2018)]{Saridakis2018tsallis}
Saridakis, E.~N., Bamba, K., Myrzakulov, R., \& Anagnostopoulos, F.~K., 2018,
\textit{Holographic dark energy through Tsallis entropy},
JCAP, 12, 012.

\bibitem[St\"{o}lzner et al.(2025)]{Stolzner2025}
St\"{o}lzner, B., et al., 2025,
\textit{KiDS-Legacy: consistency of cosmic shear measurements and joint cosmological constraints},
A\&A, arXiv:2503.19442.

\bibitem[Zurek(2003)]{Zurek2003}
Zurek, W.~H., 2003,
\textit{Decoherence, einselection, and the quantum origins of the classical},
Rev.\ Mod.\ Phys., 75, 715.

\end{thebibliography}

\end{document}
